%% kapa: exam=1 round=36 year=2025 period=1 form=A subject=real_estate subject_ko=부동산학원론 seq=8 number=88
%% source: Q-Net 감정평가사 기출문제 (cst00309) artlSeq=5231525 file=raw/1008/5231525/2025년 제36회 감정평가사 1차시험 1교시 문제지 원본.pdf
%% conversion: kordoc 4.24.51 -> tools/gen1.py (auto); review=proofread-sonnet; flags=none
\begin{question}{88}
다음 조건하에서 거미집이론(Cob-web theory)에 의한 부동산시장 A와 B의 모형 형태와 A시장과 B시장 상품의 관계는? (단, X축은 수량, Y축은 가격이고, 가격 변화에 수요는 즉각 반응하지만 공급은 시간적인 차이를 두고 반응하며, 다른 조건은 동일함)
\begin{framedtext}
○ A시장: 수요의 가격탄력성은 0.9, 공급의 가격탄력성은 1.2 \par
○ B시장: 수요곡선의 기울기는 -0.8, 공급곡선의 기울기는 0.3 \par
○ A시장의 상품가격이 4 \% 하락하면, B시장의 상품수요가 3 \% 감소함
\end{framedtext}
\medskip\noindent
\begin{tabular}{@{\hspace{0.2em}}l@{\hspace{1.2em}}l@{\hspace{1.2em}}l@{\hspace{1.6em}}l}
 & \uline{A시장} & \uline{B시장} & \uline{A와 B시장 상품의 관계} \\[0.3em]
① & 수렴형 & 순환형 & 보완재 \\
② & 수렴형 & 발산형 & 보완재 \\
③ & 발산형 & 수렴형 & 대체재 \\
④ & 발산형 & 발산형 & 대체재 \\
⑤ & 수렴형 & 수렴형 & 대체재
\end{tabular}
\end{question}
\answer{4}
